%% INTRODUCTION -- original pages xiii-xvi
\fmhead[13]{INTRODUCTION}{Introduction}{introduction}{INTRODUCTION}

The ``\textit{l}-adic representations'' considered in this book are the
algebraic analogue of the locally constant sheaves (or ``local
coefficients'') of Topology. A typical example is given by the
$\ell^n$-th division points of abelian varieties (cf.~chap.~\ref{chap:I},
\S\ref{sec:I.1.2}); the corresponding \textit{l}-adic spaces, first
introduced by Weil \cite{weil-abelian} are one of our main tools in the
study of these varieties. Even the case of dimension 1 presents non trivial
problems; some of them will be studied in chap.~\ref{chap:IV}.

The general notion of an \textit{l}-adic representation was first defined
by Taniyama \cite{taniyama} (see also the review of this paper given by
Weil in Math.~Rev., 20, 1959, rev.~1667). He showed how one can relate
\textit{l}-adic representations relative to different prime numbers
$\ell$ via the properties of the Frobenius elements (see below). In the
same paper, Taniyama also studied some abelian representations which are
closely related to complex multiplication (cf.~Weil \cite{weil-characters},
\cite{weil-cm} and Shimura--Taniyama \cite{shimura-taniyama}). These
abelian representations, together with some applications to elliptic
curves, are the subject matter of this book.

There are four Chapters, whose contents are as follows:

Chapter I begins by giving the definition and some examples of
\textit{l}-adic representations (\S\ref{sec:I.1}). In \S\ref{sec:I.2}, the
ground field is assumed to be a number field. Hence, Frobenius elements
are defined, and one has the notion of a rational \textit{l}-adic
representation: one for which their characteristic polynomials have
rational coefficients (instead of merely \textit{l}-adic ones). Two
representations corresponding to different primes are compatible if the
characteristic polynomials of their Frobenius elements are the same (at
least almost everywhere); not much is known about this notion in the non
abelian case (cf.~the list of open questions at the end of
\S\ref{sec:I.2.3}). A last section shows how one attaches L-functions to
rational \textit{l}-adic representations; the well known connection
between equidistribution and analytic properties of L-functions is
discussed in the Appendix.

Chapter II gives the construction of some abelian \textit{l}-adic
representations of a number field $K$. As indicated above, this
construction is essentially due to Shimura, Taniyama and Weil. However, I
have found it convenient to present their results in a slightly different
way, by defining first some algebraic groups over $\Q$ (the groups
$S_m$) whose representations---in the usual algebraic
sense---correspond to the sought for \textit{l}-adic representations of
$K$. The same groups had been considered before by Grothendieck in his
still conjectural theory of ``motives'' (indeed, motives are supposed to
be ``\textit{l}-adic cohomology without $l$,'' so the connection is not
surprising). The construction of these groups $S_m$ and of the
\textit{l}-adic representations attached to them, is given in
\S\ref{sec:II.2} (\S\ref{sec:II.1} contains some preliminary constructions
on algebraic groups, of a rather elementary kind). I have also briefly
indicated what relations these groups have with complex multiplication
(cf.~\S\ref{sec:II.2.8}). The last \S\ contains some more properties of
the $X(S_m)$'s.

Chapter III is concerned with the following question: let $\rho$ be an
abelian \textit{l}-adic representation of the number field $K$; can
$\rho$ be obtained by the method of chap.~\ref{chap:II}? The answer is:
this is so if and only if $\rho$ is ``locally algebraic'' in the sense
defined in \S\ref{sec:III.1}. In most applications, local algebraicity can
be checked using a result of Tate saying that it is equivalent to the
existence of a ``Hodge--Tate'' decomposition, at least when the
representation is semi-simple. The proof of this result of Tate is rather
long, and relies heavily on his theorems on $p$-divisible groups
\cite{tate-pdiv}; it is given in the Appendix.

One may also ask whether any abelian rational semi-simple \textit{l}-adic
representation of $K$ is ipso facto locally algebraic; this may well be
so, but I can prove it only when $K$ is a composite of quadratic fields;
the proof relies on a transcendency result of Siegel and Lang
(cf.~\S\ref{sec:III.3}).

Chapter IV is concerned with the \textit{l}-adic representation $V_l$
defined by an elliptic curve $E$. Its aim is to determine, as precisely as
possible, the image of the Galois group by $V_l$, or at least its Lie
algebra. Here again the ground field is assumed to be a number field (the
case of a function field has been settled by Igusa \cite{igusa}). Most of
the results have been stated in \cite{serre-lieg}, \cite{serre-resume} but
with at best some sketches of proofs. I have given here complete proofs,
granted some basic facts on elliptic curves, which are collected in
\S\ref{sec:IV.1}. The method followed is more ``global'' than the one
indicated in \cite{serre-lieg}. One starts from the fact, noticed by
Cassels and others, that the number of isomorphism classes of elliptic
curves isogenous to $E$ is finite; this is an easy consequence of
\v{S}afarevi\v{c}'s theorem (cf.~\S\ref{sec:IV.1.4}) on the finiteness
of the number of elliptic curves having good reduction outside a given
finite set of places. From this, one gets an irreducibility theorem
(cf.~\S\ref{sec:IV.2.1}). The determination of the Lie algebra of
$\Ima(V_l)$ then follows, using the properties of abelian representations
given in chap.~\ref{chap:II}, \ref{chap:III}; one has to know that $V_l$,
if abelian, is locally algebraic, but this is a consequence of the result
of Tate given in chap.~\ref{chap:III}. The variation of $\Ima(V_l)$ with
$l$ is dealt with in \S\ref{sec:IV.3}. Similar results for the local case
are given in the Appendix.
